\section{Plain-Language Protocol Summary}

\subsection{The study in one paragraph}

\draftinstr{Roughly 1,100 characters, written at the reading level of a
well-informed adult with no clinical training. Adapt the synopsis from
\texttt{inputs/\pfb phase-\pfb 1-\pfb trial-\pfb protocol.zip}, \texttt{sections/\pfb sec-\pfb 01-\pfb summary.tex}
\S1.1, and remove every term that a consent conversation would have to define
twice. Cite \texttt{Siegel2025} for the survival baseline and
\texttt{rasolute302} for the drug. State up front that this is a first-in-human
study of a combination, and that the primary question is safety, not cure.}

\subsection{What is being tested, in three parts}

\draftinstr{Three short paragraphs, roughly 500 characters each. Part one: the
drug, daraxonrasib, a RAS(ON) multi-selective inhibitor, given before and after
surgery, cite \texttt{rasolute302} and \texttt{oreilly2026darax}. Part two: the
operation, a pancreaticoduodenectomy performed with an eight-arm robotic
platform, cite \texttt{onpremwhippl} and \texttt{DutchCohort2025}. Part three:
the software, an on-premises large language model that advises and never acts,
cite \texttt{LeeAuto2024X} for where that sits on the published autonomy scale.}

\subsection{Your journey, start to finish}

\draftinstr{Roughly 900 characters introducing Figure 3, and stating explicitly
that the journey structure is adapted from the author's autonomous
single-patient journey simulation \cite{paipatientjr}, which was non-small cell
lung cancer, and that this study is PDAC. Name the four differences that matter:
the operation, the drug, the dominant early complication, and the survival
baseline. Source the distinction from
\texttt{inputs/\pfb cancer-\pfb patient-\pfb journey.zip} and \texttt{inputs/README.md}.}

\begin{pafig}[mermaid-type: flowchart LR, full page]
\node[mmgoal] (a) at (0,0) {Figure 3 slot\\journey schema\\full page, five lanes};
\node[mmin] (b) at (0,-2.2) {\ttfamily\tiny Source\\ mermaid/\\ fig-\pfb 03-\pfb journey-\pfb schema.md};
\draw[mmedge] (a) -- (b);
\end{pafig}
\vspace{-0.7cm}
\figcaption{Figure 3. The participant journey across five phases, screening through the
24-month endpoint, with every node the participant controls filled Corporate Blue\\
and every node the sponsor controls left white, so the balance of control is visible.}

\draftinstr{Replace the Figure 3 slot with the full-page mermaid-type flowchart
from \texttt{mermaid/\pfb fig-\pfb 03-\pfb journey-\pfb schema.md}. Five \texttt{mmlane} lanes at
$x = 0, 4.6, 9.2, 13.4, 18.0$; four participant \texttt{mmgoal} nodes; the
operation in \texttt{pablue1}; three dotted exits routed below all lanes at
$y = -11.2$ with right-angle routing, never diagonally through a lane.}

\subsection{Your schedule, and what it costs you in time}

\draftinstr{Roughly 800 characters introducing Figure 4 and the totals it
carries. Adapt from the Schedule of Activities in
\texttt{inputs/\pfb phase-\pfb 1-\pfb trial-\pfb protocol.zip},
\texttt{sections/\pfb sec-\pfb 01-\pfb summary.tex} \S1.3, and convert every cross in that
matrix into a statement of what the participant does. Cite
\texttt{NCITrial2024} for the general shape of trial participation.}

\begin{pafig}[d2-type: grid-rows, grid-columns]
\node[d2key] (a) at (0,0) {Figure 4 slot\\Schedule of Activities grid};
\node[d2box] (b) at (0,-1.8) {\ttfamily\tiny Source\\ d2/\\ fig-\pfb 04-\pfb soa-\pfb patient-\pfb grid.d2};
\draw[d2edge] (a) -- (b);
\end{pafig}
\vspace{-0.7cm}
\figcaption{Figure 4. The Schedule of Activities redrawn as a true grid from the
participant's chair, with each cell stating what they do rather than what is collected,\\
and a totals row giving the approximate hours of their time each visit consumes.}

\draftinstr{Replace the Figure 4 slot with the full D2-type grid from
\texttt{d2/\pfb fig-\pfb 04-\pfb soa-\pfb patient-\pfb grid.d2}: eight columns by thirteen rows of
\texttt{d2cell} rectangles sharing their rules, header row \texttt{d2cellh},
label column \texttt{d2celll}, participant-decision cells \texttt{protoblue},
activity cells \texttt{pablue1}, totals row \texttt{pagrayl}. Add the four-key
legend beneath.}

\subsection{Schedule of Activities, in the conventional form}

\draftinstr{Retain the conventional matrix as well as the grid, because a
participant may want to hand it to their own oncologist. Populate every cell
from \texttt{inputs/\pfb phase-\pfb 1-\pfb trial-\pfb protocol.zip},
\texttt{sections/\pfb sec-\pfb 01-\pfb summary.tex} \S1.3, unchanged.}

\vspace{0.30em}

\noindent\begin{tabularx}{\textwidth}{L{3.9cm} C{1.55cm} C{1.15cm} C{1.05cm} C{1.35cm} C{0.95cm} C{0.95cm} Y}
\toprule
\textbf{Activity} & \textbf{Screen -28 to -1} & \textbf{Base Day -1} & \textbf{Surg Day 0} & \textbf{Acute Day 1-7} & \textbf{Day 30} & \textbf{Day 90} & \textbf{LT q12wk} \\
\midrule
Informed consent, including Physical AI opt-out & X & & & & & & \\
Staging and KRAS G12 confirmation & X & & & & & & \\
CA 19-9 & X & X & & & X & X & X \\
ECOG performance status & X & X & & & X & X & X \\
Safety laboratory panel & X & X & X & X & X & X & X \\
Phase 0 simulation sign-off & & X & & & & & \\
USL at least 7.0 and safety matrix & & X & & & & & \\
Robotic Whipple surgery & & & X & & & & \\
Intra-operative telemetry capture & & & X & & & & \\
ISGPS fistula grading & & & & X & X & & \\
Daraxonrasib restart advisory & & & & X & & & \\
Adverse event review & X & X & X & X & X & X & X \\
RECIST imaging & X & & & & & X & X \\
PFS and OS assessment & & & & & & X & X \\
\bottomrule
\end{tabularx}

\vspace{0.30em}

\noindent Abbreviations: ISGPS, International Study Group on Pancreatic Surgery;
LT, long-term; OS, overall survival; PFS, progression-free survival; RECIST,
Response Evaluation Criteria in Solid Tumors; USL, Unified Safety Level.

\subsection{What the study cannot tell you}

\draftinstr{Roughly 700 characters, and do not soften it. A Phase 1 study of at
most eighteen participants is powered for safety and feasibility and for nothing
else. It cannot establish that this combination improves survival. Cite
\texttt{NCITrial2024} and \texttt{FDARobot2022}, the latter for the FDA's own
statement that long-term cancer outcomes have not been established for every use
of robotically assisted surgical devices. This subsection is load-bearing for
the paper's credibility and must not be trimmed in \texttt{full-\pfb patient}.}
